\chapter{晁错}

\begin{authorintro}
晁错（约前200？—前154），颍川（今河南禹州）人。少从张恢学申不害、商鞅的法家学说，以文学为太常掌故。太常寺派他到济南伏生处学《尚书》，文帝令其为太子舍人、门大夫、家令，以其辩得幸太子（即位后为景帝），太子家号曰“智囊”。数上书文帝，言削诸侯及更定法令事，文帝不听，然奇其材，迁为中大夫。景帝即位，以错为内史，错所言事辄听，宠幸倾九卿。迁为御史大夫，请上削诸侯之地。吴楚七国反，要求诛晁错以谢天下。又加窦婴、袁盎进谗言。上令晁错“衣朝衣斩东市”。晁错是位政论家，所著《论募民徙塞下书》、《论贵粟疏》等，议论犀利，分析深刻。《汉书·艺文志》法家类著录《晁错》三十一篇，已佚，今有清马国翰等人的辑本。
\end{authorintro}

\section[论贵粟疏]{论贵粟疏}

圣王\endnote{《论贵粟疏》是晁错上给汉文帝的奏疏。载《汉书·食货志上》，标题为后人所加。我国古代对农工商的政策是以农为本，以工商为末。汉文帝时代，农业和商业都获得较大发展，但农业的发展使得粮价骤降，文帝初年每石粟只卖十馀钱（见《史记·律书》），商人竞相囤积居奇。与此同时，贱五谷而贵金玉成为时尚，出现了谷贱伤农和本末倒置的新问题。晁错此疏，是为解决这些新问题而发。贵粟问题，实际是重农问题，与社会安定、边防巩固也有密切关系。疏中指出：“欲民务农，在于贵粟。贵粟之道，在于使民以粟为赏罚。”有钱的人可以纳粟买官，有罪的人可以纳粟除罪，使富人有官，农民有钱。由于作者对封建社会的市场经济与客观经济规律还不可能有正确的理解，提出的问题虽能切中时弊，但解决方案却未必切实可行。文章立论精辟，论述严谨，善用对比写法，层层深入，增强了说服力。激切的感情和流畅矫健的笔法，形成疏直峻切的风格。}\endnote{圣王：君主时代对帝王的尊称。犹“圣主”。}在上，而民不冻饥者，非能耕而食之\endnote{食（sì四）之：给人饭吃。食，作动词用。}，织而衣之也\endnote{衣之：给人衣服穿。衣，作动词用。}，为开其资财之道也\endnote{“为开”句：为民开辟生财的路子。}。故尧、禹有九年之水\endnote{尧、禹有九年之水：据《尚书·尧典》说，尧时洪水怀山襄陵，浩浩滔天，尧乃问四岳，派谁治水合适。四岳共推鲧，鲧治水九年，“绩用弗成”。又《史记·夏本纪》载：“当帝尧之时，鸿水滔天，浩浩怀山襄陵，下民其忧。……用鲧治水。九年而水不息，功用不成。”后舜摄行天子政之后，乃殛鲧于羽山，举鲧子禹治水。禹治水“居外十三年，过家门不敢入”。以上典籍均未明确说尧、禹有九年之水。又《墨子·七患》说：“《夏书》曰：‘禹七年水’。《殷书》曰：‘汤五年旱。’”但《墨子》所引《夏书》、《殷书》，不见于今本《尚书》，可能别有所据。}，汤有七年之旱\endnote{汤有七年之旱：据《淮南子·主术》载：“汤之时，七年旱，以身祷于桑林之际，而四海之云凑，千里之雨至。”上注引《墨子》谓汤有五年旱。均可见汤时确有大旱灾。}，而国亡捐瘠者\endnote{亡：同“无”。损瘠：《汉书》颜师古注：“瘠，瘦病也。言无相弃捐而瘦病者也。”}，以畜积多而备先具也\endnote{备先具：备灾之粟在灾前已准备好。}。今海内为一，土地人民之众不避汤、禹\endnote{不避：不让于，不亚于。}，加以亡天灾数年之水旱，而畜积未及者\endnote{未及者：没有达到禹、汤的水平。}，何也？地有遗利，民有馀力，生谷之土未尽垦\endnote{生谷之土：长庄稼的土地，可耕之地。垦：开垦。}，山泽之利未尽出也，游食之民未尽归农也\endnote{游食之民：无业游民。}。

民贫，则奸邪生。贫生于不足，不足生于不农\endnote{不农：不从事农业劳动。}，不农则不地著\endnote{地著（zhuó浊）：即土著（zhù住），指定居而不迁徙，使劳动力附着在土地上。}，不地著则离乡轻家，民如鸟兽\endnote{民如鸟兽：老百姓像鸟兽一样居无定所。}，虽有高城深池，严法重刑，犹不能禁也\endnote{不能禁也：指不能禁止农民的流亡与干奸邪之事。}。夫寒之于衣，不待轻暖\endnote{不待轻暖：不等待有了轻裘暖衣而后穿，即“寒不择衣”之义。}；饥之于食，不待甘旨\endnote{不待甘旨：不等待有了精美的食物而后吃。即“饥不择食”之义。}；饥寒至身，不顾廉耻。人情，一日不再食则饥，终岁不制衣则寒。夫腹饥不得食，肤寒不得衣，虽慈母不能保其子\endnote{保其子：保护、保有她的儿子。}，君安能以有其民哉\endnote{有其民：保有他的人民。}？明主知其然也，故务民于农桑\endnote{务民于农桑：尽力使人民种地养蚕。}，薄赋敛\endnote{薄赋敛：少收赋税，减轻农民负担。}，广畜积，以实仓廪\endnote{以实仓廪：以便充实粮食仓库。仓廪，储藏米谷的仓库。谷藏曰仓，米藏曰廪。一说方库曰仓，圆库曰廪。}，备水旱，故民可得而有也。

民者，在上所以牧之\endnote{上：指人君。牧：统治。}，趋利如水走下，四方亡择也\endnote{“趋利”二句：言民之追逐利益如同水向低处流，对于东西南北的方向是无所选择的。亡，同“无”。}。夫珠玉金银，饥不可食，寒不可衣，然而众贵之者，以上用之故也\endnote{以上用之故：因君主使用珠玉金银的缘故。之，指代上文的珠玉金银。}。其为物轻微易臧\endnote{其：指代珠玉金银。易臧：易藏。臧，通“藏”。}，在于把握\endnote{把握：可以一把握在手中。指珠玉金银便于携带。}，可以周海内而亡饥寒之患\endnote{周海内：即走遍天下。此句言珠玉金银的货币价值。}。此令臣轻背其主\endnote{轻背其主：轻易背叛他的主上。}，而民易去其乡，盗贼有所劝\endnote{劝：鼓励。指珠玉等对盗贼有诱惑力。}，亡逃者得轻资也\endnote{轻资：轻便的资财。指珠玉金钱。}。粟米布帛生于地，长于时，聚于力，非可一日成也。数石之重\endnote{石：此指重量单位，一石为一百二十斤。《汉书·律历志上》：“三十斤为钧，四钧为石。”}，中人弗胜\endnote{中人弗胜：中等体力的人不能胜任。}，不为奸邪所利；一日弗得而饥寒至\endnote{一日弗得：指一天得不到粟米布帛。}。是故明君贵五谷而贱金玉。

今农夫五口之家，其服役者不下二人\endnote{服役：从事于官府的劳役。}，其能耕者不过百亩。百亩之收，不过百石\endnote{石：容量单位，一石为十斗。《汉书·食货志上》：“治田百畮，岁收畮一石半。”一石约一百斤。}。春耕夏耘，秋获冬臧，伐薪樵\endnote{伐薪樵：砍伐木柴作烧柴。}，治官府\endnote{治官府：建造官府衙门。从事某种活动叫治。}，给徭役\endnote{给徭役：出劳力给官府服劳役。}；春不得避风尘，夏不得避暑热，秋不得避阴雨，冬不得避寒冻，四时之间，亡日休息\endnote{亡日：同“无日”。}。又私自送往迎来，吊死问疾，养孤长幼在其中\endnote{养孤：赡养孤老。长（zhǎng掌）幼：使幼儿成长，即培育幼儿。长，作动词用。}。勤苦如此，尚复被水旱之灾，急政暴赋\endnote{急政暴赋：用急暴的手段征取赋税。政，通“征”。}，赋敛不时，朝令而暮改当具\endnote{朝令而暮改当具：言所征之物早晨下令要晚上又改变了所应备办之物。具，备办。陈中凡先生《汉魏六朝散文选》认为此句的“改”字为衍文，即早晨下令要的东西晚上就要准备好。亦可通。}。有者半贾而卖，亡者取倍称之息\endnote{“有者”二句：言有粮的人不得不以半价贱卖，无粮的人不得不以加倍的利息以求借贷。贾，同“价”。倍称，取一偿二。}，于是有卖田宅、鬻子孙以偿责者矣\endnote{责：同“债”。}。而商贾大者积贮倍息\endnote{倍息：成倍的利息。}，小者坐列贩卖，操其奇赢\endnote{操其奇（jī机）赢：从事牟取暴利之事。奇赢，赢馀。}，日游都市，乘上之急\endnote{乘上之急：趁君王急用之时。}，所卖必倍\endnote{所卖必倍：所卖的粮食价钱必定加倍。}。故其男不耕耘，女不蚕织，衣必文采，食必粱肉；亡农夫之苦，有仟伯之得\endnote{仟伯：同“阡陌”，本指田间小路，此代指田地。}。因其富厚，交通王侯，力过吏势\endnote{力过吏势：经济实力超过官吏的势力。}，以利相倾\endnote{以利相倾：为赢利而互相倾轧。}，千里游敖\endnote{游敖：同“游遨”，游乐。作奔走周旋解亦可通。}，冠盖相望\endnote{冠盖相望：本指仕官者络绎不绝。此指商人前后不绝。冠盖，仕官者的官服和车盖。}，乘坚策肥\endnote{乘坚：乘坐坚固的车子。策肥：鞭策着肥壮的马。策，以鞭击马。}，履丝曳缟\endnote{履丝：穿着丝织品的鞋子。曳缟：拖着丝织白绢的精美衣服。}。此商人所以兼并农人，农人所以流亡者也。今法律贱商人\endnote{法律贱商人：据《史记·平准书》说：“天下已平，高祖乃令贾人不得衣丝乘车，重租税以困辱之。孝惠、高后时，为天下初定，复弛商贾之律，然市井之子孙（指商人的子孙）亦不得仕官为吏。”此可见汉初之法律是贱视商人的。}，商人已富贵矣；尊农夫，农夫已贫贱矣。故俗之所贵\endnote{俗之所贵：指农人。当时四民的排位是士、农、工、商。这就体现出世人的贵农人而贱商人。}，主之所贱也；吏之所卑\endnote{吏之所卑：官吏所卑视的。实指农民。}，法之所尊也\endnote{法之所尊：法律所尊重的。亦指农民。}。上下相反，好恶乖迕\endnote{乖迕（wǔ午）：乖违，矛盾。}，而欲国富法立，不可得也。

方今之务\endnote{务：致力、从事的根本。}，莫若使民务农而已矣\endnote{莫若：莫如，不如。}。欲民务农，在于贵粟。贵粟之道，在于使民以粟为赏罚\endnote{以粟为赏罚：用粮食作为赏罚的手段。}。今募天下入粟县官\endnote{募：征集，招募。入粟：纳粮。县官：即朝廷，也专指皇帝。}，得以拜爵\endnote{拜爵：授官爵。按此疏上文帝之后，文帝采纳了晁错的入粟拜爵之议，“令民入粟边，六百石爵上造（师古注：“上造，第二等爵也”），稍增至四千石为五大夫（师古注：“五大夫，第九等爵。”），万二千石为大庶长（师古注：“大庶长，第十八等爵。”），各以多少级数为差。”（《汉书·食货志上》）按：汉代爵的等第，与后代的官分九品不同，等数愈高官愈高。}，得以除罪\endnote{除罪：免罪，实际是纳粟赎罪。}。如此，富人有爵，农民有钱，粟有所渫\endnote{渫（xiè谢）：分散，疏散。}。夫能入粟以受爵，皆有馀者也。取于有馀，以供上用，则贫民之赋可损\endnote{损：减少。}，所谓损有馀，补不足，令出而民利者也。顺于民心，所补者三\endnote{补：裨益。所补者三：即可带来三个好处。}：一曰主用足\endnote{主用足：皇上的用度充足。}，二曰民赋少，三曰劝农功\endnote{劝农功：鼓励人从事农业生产。}。今令\endnote{今令：现在的法令。}：“民有车骑马一匹者\endnote{车骑：指战车战马。车骑马，似指既能驾战车又能作战马的两用马。}，复卒三人\endnote{复卒三人：当为卒者免三人，不当为卒者免三人的役钱（用《汉书》颜师古注）。复，免，除。}。”车骑者，天下武备也，故为复卒。神农之教曰\endnote{神农：上古传说中的帝王，始教民为耒耜，兴农业，故称神农氏。《汉书·艺文志》农家著录神农二十篇，注云是六国时（战国时）诸子疾时怠于农事，托之神农。颜师古注引刘向《别录》说：“疑李悝及商君所说。”可见《神农》是伪书。今已佚。}：“有石城十仞\endnote{石城：用石筑起的城墙。喻城之坚固，犹“金城”。仞：长度单位，一仞为周尺八尺，汉尺七尺。}，汤池百步\endnote{汤池：以沸汤为池，不能接近。喻其严固。百步：指汤池之宽。}，带甲百万，而亡粟，弗能守也。”以是观之，粟者，王者大用\endnote{王者大用：对统治天下的帝王用处最大。}，政之本务\endnote{政之本务：治政的根本急务。}。令民入粟受爵，至五大夫以上\endnote{五大夫：爵位名。春秋时始置，至秦代为二十等爵的第九级，地位在左庶长之下。汉代与秦同。}，乃复一人耳\endnote{复一人：即复卒一人。见注〔77〕“复卒三人”注。}，此其与骑马之功相去远矣\endnote{骑马之功：指上文所言出“车骑马一匹者”的功劳，复卒三人。}。爵者，上之所擅\endnote{擅：专有。}，出于口而亡穷\endnote{“出于口”句：言皇帝一开口便可封官可不受名额限制。按：滥加封爵，亦不可取。}；粟者，民之所种，生于地而不乏。夫得高爵与免罪，人之所甚欲也。使天下人入粟于边，以受爵免罪，不过三岁，塞下之粟必多矣\endnote{塞下之粟：指供边塞用的军粮。}。

\printendnotes